% TOP_PROBLEM: {"id": 2307075, "title": "Research Problems in Function Theory — Problem 7.75", "category_id": 8, "category": {"id": 8, "name": "analysis", "display_name": "Analysis", "description": "Limits, continuity, calculus, and function theory.", "slug": "analysis", "order_index": 8, "created_at": "2026-07-31T15:26:25.671Z"}, "status": "partially_solved"}
% FIRST_POSTED: null
% ENTRY_KIND: statement_only
% Imported from: batch06.tex; UnsolvedMath ID 2307075
% Compile twice: lualatex 2307075.tex
\documentclass[11pt]{article}
\usepackage{fontspec}
\setmainfont{Latin Modern Roman}
\usepackage{amsmath,amssymb,mathtools,unicode-math}
\setmathfont{Latin Modern Math}
\usepackage{xurl,hyperref}
\hypersetup{hidelinks,pdftitle={UnsolvedMath Problem 2307075}}
\newcommand{\sref}[2]{\hyperlink{source:#1:#2}{[S#2]}}
\newcommand{\cataloguescope}[1]{\paragraph{Mathematical question.}#1}
\begin{document}
% UnsolvedMath ID 2307075; source catalogue code AMR-022-7075
\setcounter{section}{2307074}
\section{Research Problems in Function Theory — Problem 7.75}\label{2307075}
{\small\sffamily UnsolvedMath ID 2307075 · Analysis}
\subsection{Definitions and mathematical statement}

\paragraph{Shared notation.}$\mathbb N=\{1,2,\ldots\}$, and $\mathbb Z,\mathbb Q,\mathbb R,\mathbb C$ denote the integers, rationals, reals and complex numbers. Any other conventions are specified locally.


For compact $E\subset\mathbb C$ define its analytic capacity
\[\gamma(E)=\sup\{|f'(\infty)|: f\in\mathcal O(\widehat{\mathbb C}\setminus E),\ |f|\le1,\ f(\infty)=0\},
\qquad f'(\infty)=\lim_{z\to\infty}zf(z).
\]
For an arbitrary bounded set $A$, its inner analytic capacity is $\gamma(A)=\sup\{\gamma(L):L\subset A,\ L\text{ compact}\}$.Identify $\mathbb C$ with $\mathbb R^2$. A $C^1$ diffeomorphism has a $C^1$ inverse; $\operatorname{GL}(2,\mathbb R)$ consists of invertible real linear maps. The four-corner quarter-square Cantor set is $C\times C$, where $C=\{\sum_{j\ge1}3\varepsilon_j4^{-j}:\varepsilon_j\in\{0,1\}\}$.
\paragraph{Statement recorded in the supplied source.}
Prove or disprove each assertion: (a) $\gamma(E)=0$ if and only if $\gamma(\phi(E))=0$ for every compact $E$ and every $C^1$ diffeomorphism $\phi:\mathbb C\to\mathbb C$; (b) the same equivalence for every $\phi\in\operatorname{GL}(2,\mathbb R)$; (c) a universal $K>0$ satisfies $\gamma(E\cup F)\le K(\gamma(E)+\gamma(F))$ for all compact $E,F$, and in particular can $K=1$ be taken? (d) If compact $K$ has $\gamma(K)=0$, must $\gamma(E\setminus K)=\gamma(E)$ for every compact $E$, with the left side interpreted as inner capacity? Include the four-corner Cantor case of (d). The original update records Tolsa's proof of (c) for some universal constant, without resolving the proposed value $1$. \sref{2307075}{2}
\subsection{Short English statement}
Study preservation of zero analytic capacity by smooth or linear changes of coordinates, universal subadditivity and its sharp constant, and removal of zero-capacity compact sets.
\subsection{Sources}
\begin{description}
\item[\hypertarget{source:2307075:1}{S1}] ULAM.AI and the UnsolvedMath Contributors, UnsolvedMath: A Curated Collection of Open Mathematics Problems, record 2307075 (AMR-022-7075). CC BY 4.0. \url{https://huggingface.co/datasets/ulamai/UnsolvedMath}.
\item[\hypertarget{source:2307075:2}{S2}] Walter K. Hayman and Edward F. Lingham, Research Problems in Function Theory (2018), Problem 7.75, its complete update and relevant preceding definitions. \url{https://arxiv.org/abs/1809.07200}.
\end{description}
\label{end:2307075}
\end{document}
